\chapter{Eliminacja Gaussa i układy równań liniowych}
\label{ch:lin5-eliminacja-gaussa}
\index{eliminacja Gaussa@Eliminacja Gaussa}
\index{operacje elementarne@Operacje elementarne na wierszach}
\index{postac schodkowa@Postać schodkowa macierzy}

\noindent\textbf{Cel rozdziału.}
W Rozdziale~\ref{ch:lin4-macierze-zmiana-bazy}
opisaliśmy układy liniowe przez równania \(Ax=b\),
ich rozwiązania i rząd macierzy.
Obecnie przedstawimy skończony algorytm, który
pozwala obliczać rząd, rozstrzygać rozwiązalność,
wyznaczać wszystkie rozwiązania oraz znajdować
bazy jądra i obrazu. Wyprowadzimy go z operacji
na wierszach i udowodnimy poprawność, zamiast
traktować przekształcenia macierzy jako przepisy
rachunkowe bez uzasadnienia.

W całym rozdziale \(K\) jest ciałem z \(0_K\ne1_K\).
Wszystkie macierze i przestrzenie są nad \(K\),
a \(m,n\in\mathbb N_0\), o ile nie wskazano inaczej.
Nie zakładamy uporządkowania ciała ani możliwości
dzielenia przez zero. Przypadki pustych macierzy
interpretujemy zgodnie z
Definicją~\ref{def:lin4-macierz}.

\section{Elementarne operacje na wierszach}
\label{sec:lin5-operacje-wierszowe}

\begin{definicja}[Trzy typy operacji elementarnych]
\label{def:lin5-operacje}
Niech \(m,n\in\mathbb N_0\) i
\(A\in\operatorname{Mat}_{m\times n}(K)\).
\emph{Elementarną operacją na wierszach} nazywamy
jedno z następujących przekształceń:
\begin{enumerate}[label=(\roman*)]
\item zamianę miejscami dwóch różnych wierszy
\(R_i\leftrightarrow R_j\), \(i\ne j\);
\item pomnożenie wiersza \(R_i\) przez skalar
\(\lambda\in K\setminus\{0_K\}\),
czyli \(R_i\leftarrow\lambda R_i\);
\item dodanie do \(R_i\) wielokrotności innego wiersza,
czyli \(R_i\leftarrow R_i+\lambda R_j\), gdzie
\(i\ne j\) i \(\lambda\in K\).
\end{enumerate}
W każdym przypadku indeksy wierszy należą do
\(\{1,\ldots,m\}\). Dla \(m=0\) nie ma
żadnej operacji do wykonania.
Dwie macierze tego samego rozmiaru nazywamy
\emph{wierszowo równoważnymi}, gdy jedną można
otrzymać z drugiej przez skończony ciąg
operacji elementarnych, także ciąg pusty.
\end{definicja}

\begin{stwierdzenie}[Odwracalność operacji elementarnych]
\label{prop:lin5-odwracalne-operacje}
Niech \(m,n\in\mathbb N_0\) i
\(A\in\operatorname{Mat}_{m\times n}(K)\).
Każda operacja elementarna ma operację odwrotną
tego samego typu. Istnieje ponadto
odwracalna macierz \(E\in\operatorname{Mat}_{m\times m}(K)\)
taka, że wynik operacji na \(A\) jest równy \(EA\).
Dla dowolnych wierszowo równoważnych macierzy
\(A,B\in\operatorname{Mat}_{m\times n}(K)\)
istnieje odwracalna macierz
\(P\in\operatorname{Mat}_{m\times m}(K)\)
spełniająca \(B=PA\).
\end{stwierdzenie}
\begin{proof}
Zamiana wierszy odwraca się przez tę samą zamianę.
Mnożenie przez \(\lambda\ne0_K\) odwraca się
mnożeniem przez \(\lambda^{-1}\), a dodanie
\(\lambda R_j\) do \(R_i\) odwraca się dodaniem
\(-\lambda R_j\). Warunek \(i\ne j\) zapewnia,
że w ostatnim przypadku wiersz \(R_j\)
pozostaje niezmieniony.

Niech \(E\) powstaje z \(I_m\) przez tę
samą operację. Z definicji mnożenia macierzy
\(i\)-ty wiersz \(EA\) jest kombinacją wierszy
\(A\) o współczynnikach z \(i\)-tego wiersza \(E\).
Dlatego \(EA\) jest wynikiem żądanej operacji.
Macierz odpowiadająca operacji odwrotnej
jest dwustronną odwrotnością \(E\), gdyż
obie operacje zastosowane kolejno do \(I_m\)
dają \(I_m\). Zatem \(E\) jest odwracalna
na mocy Definicji~\ref{def:lin4-odwracalna}.

Dla ciągu operacji odpowiadających macierzom
\(E_1,\ldots,E_s\) otrzymujemy
\(B=E_s\cdots E_1 A\).
Iloczyn \(P=E_s\cdots E_1\) jest odwracalny
na mocy Stwierdzenia~\ref{prop:lin4-iloczyn-odwrotnych}.
Przy pustym ciągu bierzemy \(P=I_m\).
Rachunek działa również dla pustych rozmiarów.
\end{proof}

\begin{wniosek}[Równoważne układy liniowe]
\label{cor:lin5-rownowazne-uklady}
Niech \(A\in\operatorname{Mat}_{m\times n}(K)\),
\(b\in K^m\), a \(P\) będzie odwracalną
macierzą rozmiaru \(m\times m\).
Wtedy układy \(Ax=b\) i \(PAx=Pb\)
mają dokładnie ten sam zbiór rozwiązań.
W szczególności można wykonywać te same
operacje elementarne na wierszach macierzy
\(A\) i odpowiednich wyrazach kolumny \(b\).
\end{wniosek}
\begin{proof}
Dla każdego \(x\in K^n\) z \(Ax=b\)
po pomnożeniu przez \(P\) otrzymujemy
\(PAx=Pb\). Odwrotnie, pomnożenie
\(PAx=Pb\) przez \(P^{-1}\) daje \(Ax=b\).
Taką samą operację na \(A\) i \(b\)
realizuje mnożenie \([A\mid b]\) przez
macierz elementarną z lewej strony.
\end{proof}

\begin{uwaga}[Czego nie zachowują operacje wierszowe?]
\label{rem:lin5-obraz-wiersze}
Jeżeli \(B=PA\) i \(P\) jest odwracalna,
to \(\ker B=\ker A\), a
\(\operatorname{im}B=P(\operatorname{im}A)\).
W szczególności \(\operatorname{rk}B=\operatorname{rk}A\),
ale obrazy \(A\) i \(B\) nie muszą być
identycznymi podprzestrzeniami \(K^m\).
Na przykład dla
\(A=(1,0)^{\mathsf T}\) i \(B=(0,1)^{\mathsf T}\)
zamiana wierszy zachowuje rząd \(1\),
lecz zmienia obraz.
\end{uwaga}

\section{Postać schodkowa i algorytm eliminacji}
\label{sec:lin5-postac-schodkowa}

\begin{definicja}[Postać schodkowa]
\label{def:lin5-schodkowa}
Niech \(R\in\operatorname{Mat}_{m\times n}(K)\).
\emph{Pierwszym niezerowym wyrazem}
niezerowego wiersza nazywamy jego wyraz
o najmniejszym indeksie kolumny, w której
występuje skalar różny od \(0_K\).
Macierz \(R\) jest w \emph{postaci schodkowej},
gdy:
\begin{enumerate}[label=(\roman*)]
\item wszystkie wiersze zerowe leżą poniżej niezerowych;
\item indeksy kolumn pierwszych niezerowych wyrazów
kolejnych niezerowych wierszy ściśle rosną.
\end{enumerate}
Pierwsze niezerowe wyrazy nazywamy
\emph{elementami wiodącymi}, a ich kolumny
\emph{kolumnami bazowymi} (kolumnami z pivotem).
Postać schodkowa jest \emph{zredukowana},
gdy dodatkowo każdy element wiodący jest
równy \(1_K\) i stanowi jedyny niezerowy wyraz
w swojej kolumnie.
Dla \(m=0\) lub \(n=0\) jedyna macierz
ma postać schodkową zredukowaną.
\end{definicja}

\begin{twierdzenie}[Istnienie postaci schodkowej i zredukowanej]
\label{thm:lin5-gauss}
Dla dowolnych \(m,n\in\mathbb N_0\)
i \(A\in\operatorname{Mat}_{m\times n}(K)\)
istnieją wierszowo równoważne z \(A\)
macierze \(U,R\) tego samego rozmiaru,
z których \(U\) jest schodkowa,
a \(R\) schodkowa zredukowana.
Można je otrzymać skończoną liczbą
elementarnych operacji na wierszach.
\end{twierdzenie}
\begin{proof}
Podamy algorytm wraz z uzasadnieniem.
Ustawiamy numer aktualnego wiersza \(r=1\)
i przeglądamy kolumny \(j=1,\ldots,n\)
w rosnącej kolejności. Jeżeli \(r>m\),
kończymy; dalsze kolumny nie mogą
dostarczyć nowych wierszy niezerowych.

W rozpatrywanej kolumnie \(j\) szukamy
wiersza \(k\geq r\) z niezerowym wyrazem
w tej kolumnie. Gdy takiego wiersza nie ma,
przechodzimy do następnej kolumny.
Gdy istnieje, zamieniamy wiersze \(k\) i \(r\)
(jeżeli \(k=r\), zamiana nie jest potrzebna).
Element \(a_{rj}\) jest teraz niezerowy,
więc wolno pomnożyć wiersz \(r\)
przez \(a_{rj}^{-1}\), uzyskując jedynkę.
Od każdego wiersza \(i>r\) odejmujemy
\(a_{ij}\) razy nowy wiersz \(r\);
zerujemy zatem wszystkie wyrazy pod pivotem.
Zwiększamy \(r\) o \(1\).

Po wybraniu pivotu w kolumnie \(j\)
wszystkie wiersze poniżej \(r\) mają zera
w kolumnach wcześniej wybranych:
były tam wyzerowane podczas poprzednich
etapów, a kombinacje tych wierszy
nie wprowadzają nowych niezerowych wyrazów.
Jeżeli pewną kolumnę pominięto z powodu
samych zer pod aktualnym wierszem,
kolejne operacje na tych niższych wierszach
również zachowują w tej kolumnie zera.
Stąd kolejne wybrane pivoty występują
w ściśle rosnących kolumnach,
a wiersze niewybrane, jeśli pozostały,
są zerowe. Otrzymana macierz \(U\)
ma postać schodkową. Algorytm przegląda
dokładnie \(n\) kolumn i wybiera nie więcej
niż \(\min\{m,n\}\) pivotów, zatem się kończy.

Aby uzyskać postać zredukowaną, przechodzimy
po pivotach od ostatniego ku pierwszemu.
W każdym wierszu wiodący wyraz wynosi już \(1_K\).
Za jego pomocą odejmujemy odpowiednie
wielokrotności tego wiersza od wszystkich
wierszy położonych wyżej, zerując wyrazy
nad pivotem. Ponieważ wiersz z późniejszym
pivotem ma zera we wszystkich wcześniejszych
kolumnach bazowych, operacje te nie niszczą
zer uzyskanych w poprzednich krokach
ani samych jedynek wiodących.
Po zakończeniu każda kolumna bazowa ma
jedną jedynkę i same zera. Powstała
macierz \(R\) jest schodkowa zredukowana
i wierszowo równoważna z \(A\).

Jeżeli jeden z rozmiarów jest zerowy,
nie wykonujemy żadnego kroku;
jedyna macierz już jest zredukowana.
\end{proof}

\begin{przyklad}[Dwa etapy eliminacji]
\label{prz:lin5-pierwsza-eliminacja}
Nad \(\mathbb R\) niech
\[
 A=\begin{pmatrix}
 1&2&1&0\\
 2&4&3&1\\
 1&2&2&1
 \end{pmatrix}.
\]
Odejmujemy dwukrotność pierwszego wiersza
od drugiego i pierwszy wiersz od trzeciego:
\[
 A\ \sim\
 \begin{pmatrix}
 1&2&1&0\\
 0&0&1&1\\
 0&0&1&1
 \end{pmatrix}
 \ \sim\
 U=\begin{pmatrix}
 1&2&1&0\\
 0&0&1&1\\
 0&0&0&0
 \end{pmatrix}.
\]
Elementy wiodące są w kolumnach \(1\) i \(3\).
Odejmując drugi wiersz od pierwszego,
otrzymujemy postać zredukowaną
\[
 R=\begin{pmatrix}
 1&2&0&-1\\
 0&0&1&1\\
 0&0&0&0
 \end{pmatrix}.
\]
Każda strzałka jest ciągiem odwracalnych
operacji; nie wykonywaliśmy operacji
na kolumnach.
\end{przyklad}

\section{Liczba pivotów i rząd macierzy}
\label{sec:lin5-rzad}

\begin{twierdzenie}[Rząd macierzy schodkowej]
\label{thm:lin5-rzad-pivoty}
Niech \(m,n\in\mathbb N_0\),
\(A\in\operatorname{Mat}_{m\times n}(K)\),
a \(U\) będzie dowolną postacią schodkową
wierszowo równoważną z \(A\).
Jeżeli \(U\) ma \(r\) niezerowych wierszy,
to
\[
 \operatorname{rk}A=r,\qquad
 \dim_K\ker A=n-r.
\]
W szczególności liczba pivotów nie zależy
od wyborów wykonanych w algorytmie.
\end{twierdzenie}
\begin{proof}
Niech \(p_1<\cdots<p_r\) będą indeksami
kolumn elementów wiodących.
W układzie jednorodnym \(Ux=0\)
pozostałe \(n-r\) współrzędnych, tj.
\(x_j\) dla \(j\notin\{p_1,\ldots,p_r\}\),
można wybrać dowolnie. Rzeczywiście,
zaczynając od ostatniego niezerowego wiersza,
wyznaczamy jednoznacznie \(x_{p_r}\),
następnie \(x_{p_{r-1}}\), aż do \(x_{p_1}\).
Jest to możliwe, ponieważ każdy
element wiodący jest niezerowy
i ma odwrotność w ciele \(K\).

Rzut z \(\ker U\) na kolumnę \(n-r\)
współrzędnych swobodnych jest liniową
bijekcją na \(K^{n-r}\):
istnienie i jednoznaczność wynika
z opisanego podstawiania wstecz,
a liniowość z liniowości rzutu.
Zatem \(\dim_K\ker U=n-r\).
Operacje wierszowe nie zmieniają jądra,
więc \(\ker A=\ker U\).
Twierdzenie~\ref{thm:lin4-rzad-kolumn}
daje
\(n=\dim_K\ker A+\operatorname{rk}A\),
skąd \(\operatorname{rk}A=r\).
Gdy \(r=0\) albo \(n=0\), argument
oznacza po prostu wybór wszystkich
współrzędnych jako swobodnych lub
pusty układ współrzędnych.
\end{proof}

\begin{wniosek}[Kryteria liczby pivotów]
\label{cor:lin5-pivoty-odwracalnosc}
Niech \(A\in\operatorname{Mat}_{m\times n}(K)\),
gdzie \(m,n\in\mathbb N_0\).
Macierz \(A\) wyznacza iniektywne
\(L_A\colon K^n\to K^m\) wtedy i tylko wtedy,
gdy każda kolumna jest bazowa w jej postaci
schodkowej. Wyznacza surjekcję wtedy i tylko
wtedy, gdy liczba pivotów wynosi \(m\).
Jeżeli \(m=n\), jest odwracalna wtedy i
tylko wtedy, gdy jej postacią schodkową
zredukowaną jest \(I_n\).
\end{wniosek}
\begin{proof}
Iniektywność jest równoważna
\(\dim\ker A=0\), czyli \(r=n\).
Surjektywność jest równoważna
\(\operatorname{rk}A=m\), czyli \(r=m\).
Gdy \(m=n\), odwracalność jest
równoważna \(r=n\) na mocy
Twierdzenia~\ref{thm:lin4-odwracalnosc-kryteria}.
Zredukowana macierz kwadratowa z \(n\)
pivotami musi mieć pivot w każdej
kolumnie i wierszu, a więc być równa \(I_n\).
Implikacja odwrotna wynika z liczby pivotów.
\end{proof}

\section{Rozwiązywanie układów i kryterium zgodności}
\label{sec:lin5-uklady}

\begin{definicja}[Macierz rozszerzona]
\label{def:lin5-rozszerzona}
Dla \(A\in\operatorname{Mat}_{m\times n}(K)\)
i \(b\in K^m\) przez \([A\mid b]\)
oznaczamy macierz rozmiaru \(m\times(n+1)\),
której pierwsze \(n\) kolumn tworzą \(A\),
a ostatnią kolumną jest \(b\).
Operacje wierszowe wykonujemy na całej
macierzy \([A\mid b]\), a nie jedynie na \(A\).
\end{definicja}

\begin{twierdzenie}[Kryterium zgodności układu]
\label{thm:lin5-zgodnosc}
Niech \(m,n\in\mathbb N_0\),
\(A\in\operatorname{Mat}_{m\times n}(K)\)
oraz \(b\in K^m\).
Następujące warunki są równoważne:
\begin{enumerate}[label=(\roman*)]
\item układ \(Ax=b\) ma rozwiązanie;
\item \(b\in\operatorname{im}A\);
\item \(\operatorname{rk}A=\operatorname{rk}[A\mid b]\);
\item po sprowadzeniu \([A\mid b]\) do postaci
schodkowej nie występuje wiersz
\((0,\ldots,0\mid\beta)\) z \(\beta\ne0_K\).
\end{enumerate}
Jeżeli układ jest zgodny i
\(r=\operatorname{rk}A\), to jego zbiór rozwiązań
jest warstwą podprzestrzeni wymiaru \(n-r\).
\end{twierdzenie}
\begin{proof}
Równoważność (i) i (ii) jest
Twierdzeniem~\ref{thm:lin4-zbior-rozwiazan}.
Obraz macierzy \([A\mid b]\) jest
przestrzenią rozpiętą przez wszystkie
kolumny \(A\) i dodatkową kolumnę \(b\).
Ma ten sam wymiar co obraz \(A\)
dokładnie wtedy, gdy \(b\) już należy
do przestrzeni rozpiętej przez kolumny \(A\):
jeśli nie należy, dołączenie \(b\)
zwiększa wymiar o jeden.
To dowodzi równoważności (ii) i (iii).

Operacje wierszowe na \([A\mid b]\)
nie zmieniają zbioru rozwiązań
na mocy Wniosku~\ref{cor:lin5-rownowazne-uklady}.
Wiersz \((0,\ldots,0\mid\beta)\),
\(\beta\ne0_K\), oznacza równanie
\(0_K=\beta\), którego żaden wektor
nie spełnia. Jeśli takiego wiersza
nie ma w postaci schodkowej, ostatnia
kolumna nie jest kolumną nowego pivotu,
a podstawianie wstecz, z dowolnym wyborem
zmiennych swobodnych, wyznacza rozwiązanie.
Równoważne są więc (i) i (iv).

Dla zgodnego układu rozwiązania mają
postać \(x_0+\ker A\) na mocy
Twierdzenia~\ref{thm:lin4-zbior-rozwiazan},
a wymiar \(\ker A\) wynosi \(n-r\)
na mocy Twierdzenia~\ref{thm:lin5-rzad-pivoty}.
Warstwa nie jest na ogół podprzestrzenią,
jeżeli \(b\ne0_{K^m}\).
\end{proof}

\begin{przyklad}[Układ zgodny i niezgodny]
\label{prz:lin5-uklad-rozszerzony}
Rozpatrzmy macierz \(A\) z
Przykładu~\ref{prz:lin5-pierwsza-eliminacja}
i \(b=(3,7,4)^{\mathsf T}\in\mathbb R^3\).
Wykonując kolejno
\(R_2\leftarrow R_2-2R_1\),
\(R_3\leftarrow R_3-R_1\),
\(R_3\leftarrow R_3-R_2\),
\(R_1\leftarrow R_1-R_2\), otrzymujemy
\[
 [A\mid b]\ \sim\
 \left(\begin{array}{rrrr|r}
 1&2&0&-1&2\\
 0&0&1& 1&1\\
 0&0&0& 0&0
 \end{array}\right).
\]
Zmienne \(x_2=s\), \(x_4=t\) są swobodne,
a równania wiodące dają
\(x_1=2-2s+t\), \(x_3=1-t\).
Stąd wszystkie rozwiązania mają postać
\[
 x=\begin{pmatrix}2\\0\\1\\0\end{pmatrix}
 +s\begin{pmatrix}-2\\1\\0\\0\end{pmatrix}
 +t\begin{pmatrix}1\\0\\-1\\1\end{pmatrix},
 \qquad s,t\in\mathbb R.
\]
Liczba zmiennych swobodnych jest równa
\(4-\operatorname{rk}A=2\).
Jeśli zamiast \(b\) weźmiemy
\(b'=(3,7,5)^{\mathsf T}\),
te same pierwsze trzy operacje wytwarzają
wiersz \((0,0,0,0\mid1)\).
Układ \(Ax=b'\) nie ma rozwiązania;
jego brak jest matematycznie różny
od nieskończenie wielu rozwiązań układu
z prawą stroną \(b\).
\end{przyklad}

\section{Bazy jądra i obrazu wyznaczane przez eliminację}
\label{sec:lin5-bazy}

\begin{twierdzenie}[Baza jądra z postaci zredukowanej]
\label{thm:lin5-baza-jadra}
Niech \(m,n\in\mathbb N_0\) i
\(A\in\operatorname{Mat}_{m\times n}(K)\).
Niech \(R\) będzie postacią schodkową
zredukowaną \(A\), a
\(p_1<\cdots<p_r\) jej kolumnami bazowymi.
Dla każdej kolumny swobodnej
\(f\notin\{p_1,\ldots,p_r\}\)
definiujemy \(v^{(f)}\in K^n\) przez warunki
\[
\begin{aligned}
 v^{(f)}_f&=1_K,\\
 v^{(f)}_{p_i}&=-R_{if}\quad(1\leq i\leq r),\\
 v^{(f)}_g&=0_K
 \quad\text{dla każdej swobodnej }g\ne f.
\end{aligned}
\]
Wówczas rodzina \(\bigl(v^{(f)}\bigr)_f\)
stanowi bazę \(\ker A\).
\end{twierdzenie}
\begin{proof}
Wiersz \(i\) macierzy \(R\) daje równanie
\[
 x_{p_i}+
 \sum_{f\ \text{swobodna}}R_{if}x_f=0_K,
\]
ponieważ w pozostałych kolumnach bazowych
ma same zera. Wektor \(v^{(f)}\)
spełnia każde z tych równań,
zatem należy do \(\ker R=\ker A\).

Dla dowolnego \(x\in\ker A\)
porównanie wszystkich współrzędnych
swobodnych i powyższe równania dają
\(x=\sum_f x_fv^{(f)}\),
czyli rodzina rozpina jądro.
Jeśli \(\sum_f\lambda_fv^{(f)}=0\),
to współrzędna swobodna o indeksie \(f\)
wynosi \(\lambda_f\), więc wszystkie
współczynniki są zerowe. Rodzina jest
liniowo niezależna. Jeśli nie ma
kolumn swobodnych, jest pusta i
\(\ker A=\{0_{K^n}\}\).
\end{proof}

\begin{twierdzenie}[Baza obrazu z kolumn oryginalnej macierzy]
\label{thm:lin5-baza-obrazu}
Niech \(A\in\operatorname{Mat}_{m\times n}(K)\),
gdzie \(m,n\in\mathbb N_0\), i niech
\(R\) będzie jej postacią schodkową
zredukowaną, z kolumnami bazowymi
o indeksach \(p_1<\cdots<p_r\).
Wówczas kolumny macierzy \emph{oryginalnej}
\(A\) o indeksach \(p_1,\ldots,p_r\)
tworzą bazę \(\operatorname{im}A\).
Nie należy zastępować ich odpowiadającymi
kolumnami \(R\), gdy szukamy bazy
obrazu \(A\) jako podprzestrzeni \(K^m\).
\end{twierdzenie}
\begin{proof}
Z odwracalności operacji elementarnych
istnieje odwracalna macierz \(P\)
taka, że \(R=PA\).
Kolumny \(R\) o indeksach \(p_i\)
są odpowiednio wektorami
\(e_i\in K^m\), \(1\leq i\leq r\).
Są więc liniowo niezależne.

Każda inna kolumna \(R\) ma niezerowe
współrzędne co najwyżej w pierwszych
\(r\) wierszach i jest kombinacją
tych kolumn bazowych:
\(R_{\bullet j}=\sum_{i=1}^r R_{ij}R_{\bullet p_i}\).
Po zastosowaniu \(P^{-1}\) otrzymujemy
\(A_{\bullet j}=\sum_{i=1}^r R_{ij}A_{\bullet p_i}\).
Odwracalność \(P\) zachowuje także
niezależność liniową wybranych kolumn.
Zatem wybrane kolumny \(A\) są bazą
przestrzeni rozpiętej przez wszystkie
kolumny, która jest równa
\(\operatorname{im}A\) na mocy
Twierdzenia~\ref{thm:lin4-rzad-kolumn}.
W przypadku \(r=0\) rodzina jest pusta
i obraz jest zerowy.
\end{proof}

\begin{przyklad}[Jądro i obraz z jednego rachunku]
\label{prz:lin5-bazy}
Dla macierzy \(A\) z
Przykładu~\ref{prz:lin5-pierwsza-eliminacja}
postać zredukowana ma kolumny bazowe
\(1,3\), a kolumny swobodne \(2,4\).
Wstawiając kolejno \((x_2,x_4)=(1,0)\)
oraz \((0,1)\), otrzymujemy bazę jądra
\[
 \left\{
 \begin{pmatrix}-2\\1\\0\\0\end{pmatrix},
 \begin{pmatrix}1\\0\\-1\\1\end{pmatrix}
 \right\}.
\]
Bazę obrazu tworzą \emph{pierwsza i trzecia
kolumna pierwotnej macierzy}:
\[
 \left\{
 \begin{pmatrix}1\\2\\1\end{pmatrix},
 \begin{pmatrix}1\\3\\2\end{pmatrix}
 \right\}.
\]
Są one niezależne, bo równanie
\(\alpha(1,2,1)+\beta(1,3,2)=0\)
daje z pierwszej współrzędnej
\(\alpha=-\beta\), a z drugiej
\(\beta=0\); zatem także \(\alpha=0\).
Wymiar jądra jest równy \(2\),
wymiar obrazu też jest równy \(2\),
a suma wynosi \(4=\dim\mathbb R^4\).
\end{przyklad}

\section{Jednoznaczność postaci zredukowanej}
\label{sec:lin5-jednoznacznosc}

\begin{twierdzenie}[Jednoznaczność postaci schodkowej zredukowanej]
\label{thm:lin5-jednoznacznosc-rref}
Niech \(m,n\in\mathbb N_0\).
Każda macierz \(A\in\operatorname{Mat}_{m\times n}(K)\)
ma dokładnie jedną postać schodkową
zredukowaną w swojej klasie
wierszowej równoważności.
W szczególności dwie macierze tego samego
rozmiaru są wierszowo równoważne wtedy
i tylko wtedy, gdy mają tę samą
postać schodkową zredukowaną.
\end{twierdzenie}
\begin{proof}
Istnienie wynika z
Twierdzenia~\ref{thm:lin5-gauss}.
Niech \(R,S\) będą dwiema zredukowanymi
postaciami macierzy \(A\).
Operacje wierszowe nie zmieniają
przestrzeni rozpiętej przez wiersze:
każdy nowy wiersz jest kombinacją starych,
a odwracalność operacji daje również
odwrotne zawieranie. Zatem \(R\) i \(S\)
mają tę samą przestrzeń wierszy \(H\subseteq K^n\).

Załóżmy, że niezerowe wiersze \(R\)
mają pierwsze niezerowe wyrazy w kolumnach
\(p_1<\cdots<p_r\). Dla niezerowej
kombinacji wierszy \(R\) pierwszy
niezerowy wyraz leży w jednej z kolumn
\(p_i\). Rzeczywiście, wybieramy
najmniejszy indeks \(i\), przy którym
współczynnik wiersza jest niezerowy.
W kolumnach wcześniejszych niż \(p_i\)
wszystkie użyte wiersze mają zera,
a wyraz w kolumnie \(p_i\) jest równy
temu współczynnikowi i jest niezerowy.
Z drugiej strony każda kolumna \(p_i\)
jest pierwszą niezerową kolumną
odpowiedniego wiersza, a zatem wszystkie
te pozycje występują w przestrzeni \(H\).
Zbiór pozycji pierwszego niezerowego wyrazu
w niezerowych wektorach \(H\) jest więc
dokładnie \(\{p_1,\ldots,p_r\}\).
Jest on określony samą przestrzenią \(H\),
dlatego macierz \(S\) ma te same
indeksy kolumn bazowych.

Niech \(u_i\) będzie \(i\)-tym niezerowym
wierszem \(R\), a \(v_i\) takim wierszem \(S\).
Ponieważ obie macierze są zredukowane,
\((u_i)_{p_j}=(v_i)_{p_j}\)
wynosi \(1_K\) dla \(j=i\)
i \(0_K\) dla \(j\ne i\).
Wektor \(v_i\in H\) jest pewną kombinacją
wierszy \(u_1,\ldots,u_r\).
Porównanie współrzędnych w kolumnach
\(p_j\) pokazuje, że współczynniki
tej kombinacji wynoszą odpowiednio
\(0_K,\ldots,1_K,\ldots,0_K\).
Zatem \(v_i=u_i\) dla każdego \(i\);
pozostałe wiersze w obu macierzach
są zerowe. Otrzymujemy \(R=S\).

Gdy \(H=\{0\}\), obie macierze są
zerowe i rozumowanie pozostaje prawdziwe.
Jeżeli dwie macierze są wierszowo
równoważne, ich postaci zredukowane
są równe z wykazanej jednoznaczności.
Odwrotnie, gdy ich postacie są równe,
odwracalność każdego kroku eliminacji
pozwala przejść od pierwszej macierzy
do drugiej przez skończony ciąg operacji.
\end{proof}

\section{Obliczanie macierzy odwrotnej}
\label{sec:lin5-odwracanie}

\begin{twierdzenie}[Metoda Gaussa--Jordana]
\label{thm:lin5-gauss-jordan-odwrotnosc}
Niech \(n\in\mathbb N_0\) i
\(A\in\operatorname{Mat}_{n\times n}(K)\).
Wykonujmy te same elementarne operacje
na obu blokach macierzy \([A\mid I_n]\),
tak aby lewy blok uzyskał postać
schodkową zredukowaną \(R\).
Wówczas \(R=I_n\) dokładnie wtedy,
gdy \(A\) jest odwracalna; w tym przypadku
prawy blok końcowej macierzy jest równy
\(A^{-1}\).
\end{twierdzenie}
\begin{proof}
Łącznie wszystkie operacje odpowiadają
mnożeniu z lewej strony przez pewną
odwracalną macierz \(P\).
Wynikiem jest więc \([PA\mid P]\),
gdzie \(R=PA\).
Jeśli \(A\) jest odwracalna,
Twierdzenie~\ref{thm:lin5-rzad-pivoty}
daje \(\operatorname{rk}R=n\);
zredukowana macierz kwadratowa rzędu \(n\)
jest równa \(I_n\).
Wtedy \(PA=I_n\) i po pomnożeniu przez
\(A^{-1}\) otrzymujemy \(P=A^{-1}\).
Odwrotnie, jeśli \(R=I_n\),
to \(PA=I_n\). Ponieważ \(P\) jest odwracalna,
mamy \(A=P^{-1}\), a więc \(A\) jest
odwracalna i \(P=A^{-1}\).
Przy \(n=0\) obie strony są pustymi
macierzami jednostkowymi i wniosek
jest nadal prawdziwy.
\end{proof}

\begin{przyklad}[Odwrotność metodą eliminacji]
\label{prz:lin5-odwrotnosc}
Niech
\[
 A=\begin{pmatrix}2&1\\1&1\end{pmatrix}
 \in\operatorname{Mat}_{2\times2}(\mathbb R).
\]
Rozszerzamy ją o \(I_2\), następnie
zamieniamy wiersze, eliminujemy,
normalizujemy pivot i zerujemy wyraz nad nim:
\[
\begin{aligned}
\left(\begin{array}{rr|rr}2&1&1&0\\1&1&0&1\end{array}\right)
&\sim
\left(\begin{array}{rr|rr}1&1&0&1\\2&1&1&0\end{array}\right)\\
&\sim
\left(\begin{array}{rr|rr}1&1&0&1\\0&-1&1&-2\end{array}\right)\\
&\sim
\left(\begin{array}{rr|rr}1&1&0&1\\0&1&-1&2\end{array}\right)\\
&\sim
\left(\begin{array}{rr|rr}1&0&1&-1\\0&1&-1&2\end{array}\right).
\end{aligned}
\]
Zatem
\[
 A^{-1}=\begin{pmatrix}1&-1\\-1&2\end{pmatrix}.
\]
Wynik zgadza się z
Przykładem~\ref{prz:lin4-odwrotnosc-2}.
Nie wyznaczaliśmy odwrotności
przez zgadywanie ani przez wyznacznik.
\end{przyklad}

\section{Podsumowanie}
\label{sec:lin5-podsumowanie}

Operacje elementarne na wierszach są
odwracalne i odpowiadają mnożeniu
z lewej strony przez macierze odwracalne.
Eliminacja Gaussa przekształca każdą macierz
w postać schodkową, a eliminacja
Gaussa--Jordana w jednoznacznie wyznaczoną
postać schodkową zredukowaną.
Liczba pivotów jest rzędem macierzy,
a liczba zmiennych swobodnych
w układzie jednorodnym wynosi
\(n-\operatorname{rk}A\).
Układ \(Ax=b\) jest zgodny dokładnie wtedy,
gdy \(\operatorname{rk}A=\operatorname{rk}[A\mid b]\).
Z postaci zredukowanej można obliczyć
bazę jądra oraz wskazać indeksy kolumn
pierwotnej macierzy tworzących bazę obrazu.
W przypadku macierzy kwadratowej
redukcja \([A\mid I_n]\) pozwala
rozstrzygnąć odwracalność i znaleźć \(A^{-1}\).
